% ======================================================================
% Capítulo: Estudio comparativo de estrategias de \emph{chunking}
% Fase 4 del anteproyecto. Extensión orientativa: 5--7 páginas. Resultados completos en el apéndice~\ref{ap:chunking}.
% ======================================================================
\chapter{Estudio comparativo de estrategias de \emph{chunking}}
\label{cap:estudio-chunking}

\section{Estrategias comparadas}
\label{sec:estrategias}
% TODO: Fixed-size con solapamiento, recursive, semantic y jerárquica padre/hijo, con sus parámetros.

\section{Diseño experimental}
\label{sec:diseno-experimental}
% TODO: Corpus fijo, preguntas de referencia con \emph{ground truth}, mismos modelos de embeddings y generación.

\section{Entorno de evaluación local}
\label{sec:entorno-local}
% TODO: RAGAS en local con Ollama como juez; justificación por coste.

\section{Métricas}
\label{sec:metricas}
% TODO: Precisión y recall de contexto, fidelidad, relevancia, tasa de alucinación, latencia y coste de indexación.

\section{Resultados}
\label{sec:resultados}
% TODO: Tablas y gráficas por estrategia y métrica.

\section{Discusión}
\label{sec:discusion}
% TODO: Qué estrategia gana, con qué coste, y si confirma la jerárquica en producción.

\section{Amenazas a la validez}
\label{sec:validez}
% TODO: Tamaño del conjunto de preguntas, modelo local como juez, representatividad del corpus.
